\subsection{ZO0K' variant 2 (legacy)}
\label{s:zook:variant:2}

Given WHIR with $1\leq r\leq n$ rounds, consisting of $k_1,\ldots,k_r$ sumcheck rounds,
and $n = k_1 + k_2 + \ldots + k_r$.
Let $n_0 = 0$, $n_1 = k_1$, $n_2 = k_1 + k_2 $, \ldots

Instead of adding the masks $\msk_1,\ldots,\msk_{r-1}\in F[X]^{<2^m}$ stepwise, we overlay them altogether with the input polynomial $p_0(X_1,\ldots, X_n)$ during its commit procedure, using the joint mask 
from the sparse univariate
\begin{align}
\label{e:joint:mask}
\msk(X_1,\ldots,X_n) \simeq  \msk_0(X) + \msk_1(X^{2^{n_1}}) + \ldots + \msk_{r-1}(X^{2^{n_r}}).
\end{align}
However, unlike in the zkWHIR 2.0 proposal, we also use quadratic sumcheck
blinders $h_1,\ldots h_n$, which we mix into the sumcheck at the very beginning.

\subsubsection{Masked commitment.}
Given the multilinear $p_1(X_1,\ldots, X_n)$ we commit to it in the following manner.
\begin{itemize}
\item 
Sample masks and mask blinders $\msk_i, b_i \sample F[X]^{2^m}$, $i=1,\ldots, r$, and commit to
\[
\hat\msk_i(X) = \msk_i(X) + X^{2^m} \cdot b_i(X),
\]
% $i=0,\ldots r-1$, 
% \[
% \hat \msk_i(X_1,\ldots, X_m) = \msk(X_1,\ldots, X_m) + b_i(X_1,\ldots, X_m), 
% \]
$i=0,\ldots, r-1$, as a codewords from $C'$.

\item 
Commit
\begin{align*}
\hat p_0(X_1,\ldots, X_n) = p_0(X_1,\ldots, X_n) + \msk(X_1,\ldots, X_n),
%+ \sum_{i=0}^{r-1} \msk_i(X_{n_i + 1}, \ldots ,X_{n_i + m}),
\end{align*}
as codeword from $C_n$, where $\msk$ is the composed masked from \eqref{e:joint:mask}.
Notice that the univariate in \eqref{e:joint:mask} consists only of $r\times 2^m$ coefficients, thus adding them to $p_0$ is a ``local'' modification.
\end{itemize}

\subsubsection{The protocol.}
Given $p_1(X_1,\ldots, X_n)$, committed as above, we prove the inner product claim 
\[
\sum_{\vec x \in H_n} p_0(\vec x)\cdot K(\vec x) = c_0
\]
as follows:

\begin{enumerate}
\item 
The prover samples the sumcheck blinder $h_j(X)\sample F[X]^{<3}$, and commits them as codewords from $C'$, via
\[
\hat h_j(X) = h_j(X) + X^{2^m} \cdot c_j(X),
\]
% $j=1,\ldots,n$, 
with $c_j(X)\sample F[X]^{2^m}$.
The sumcheck mask is the multivariate 
\[
h(X_1,\ldots, X_n) = h_1(X_1) + \ldots + h_n(X_n).
\]
Note that with $\hat h_1,\ldots, \hat h_n$, we have implicitly committed to the hypercube sum
\[
\langle 1, h\rangle_{H_n} = \sum_{\vec x \in\{0,1\}^n} h_1(x_1)+ \ldots + h_n(x_n)
= \sum_{j=1}^n h_j(0) + h_j(1),
\]
and the prover does not need to announce it explicitly. 

\item
The verifier answers with $\gamma\sample F$, and both the prover and the verifier now run WHIR on $\hat p_0$, $\msk$, and $h$ on the joint claim
\begin{align*}
\sum_{\vec x\in H_n} 
\left(\hat p_0(\vec x) - \msk(\vec x) \right) 
\cdot K(\vec x) + \gamma\cdot h(\vec x)  &= c_0 + \gamma \cdot \langle 1, h\rangle_{H_n}.
\end{align*}
After the final round of WHIR the claim is reduced to 
\begin{align}
\label{e:final:claim}
\left(\hat p_0(\vec \lambda) - \msk(\vec \lambda) \right) 
\cdot K'(\vec \lambda) + \gamma\cdot h(\vec \lambda)  &= c_0 + \gamma \cdot \langle 1, h\rangle_{H_n}, %= s_0,
\end{align}
where $\vec\lambda\in F^n$ are the folding challenges of WHIR, and $K'$ is the modified kernel, which has all the single-point queries absorbed. 

\item 
The final claim \eqref{e:final:claim} is linear in the univariates $\msk_0,\ldots, \msk_{r-1}$ and $h_1,\ldots, h_n$, hence also in their committed versions 
\[
\hat\msk_0,\ldots, \hat\msk_{r-1},\hat h_1,\ldots, \hat h_n \in F[X]^{<2^{m+1}}.
\]
The prover commits to the random polynomials
\[
m_0,\ldots, m_{r-1}, m_r, \ldots, m_{r+n-1} \sample F[X]^{<2^{m+1}} 
\]
as codewords from $C'$, reveals the value $v$ of their linear functional.
Both claims are batched using another verifier randomness, and the combined claim is proven by complete reveal of its witness, together with the collinearity checks at the query points. 
\end{enumerate}


\begin{lem}
Let $Q_0$ and  $Q_1, \ldots, Q_{r-1}$ be any choice of query sets for $\hat f_0$ and its foldings $\hat f_1$, \ldots, $\hat f_{r-1}$ committed in WHIR.
(That is, we fix all verifier challenges in WHIR.)
Then the distribution of 
\[
\left(\hat f_0|_{Q_0}, \ldots, \hat f_{r-1}|_{Q_{r-1}}\right)
\in F^{|Q_0|+\ldots +|Q_{r-1}|}
\]
is uniform on $F^{|Q_0|+\ldots +|Q_{r-1}|}$, when varying over all masks $\msk_0,\ldots,\msk_{r-1}\in F[X]^{<2^m}$.
(Or course, independent of $h_1,\ldots, h_n$.)
% ($\hat f_r$ is revealed in full.)
\end{lem}
\begin{proof}
I think best by induction, "bottom-up".
Purely univariate.
$\Box$
\end{proof}


\begin{lem}
Again for $Q_0, Q_1,\ldots, Q_{r-1}$ and the other verifier challenges in WHIR fixed.
The distribution of the sumcheck polynomials 
$q_1(X),\ldots, q_n(X)$ is uniform on 
\[
\mathcal R_k = 
\left\{
(q_{n_k+1},\ldots, q_{n_{k+1}})\in \left(F[X]^{<3}\right)^n \::\:
\begin{matrix}
c = q_1(0) + q_1(0)
\\
q_i(\lambda_i) = q_{i+1}(0) + q_{i+1}(1),
\\
\text{for }i=1,\ldots,n-1,
\\
q_n(\lambda_n) = \text{last sumcheck claim} 
\\
\text{with the (todo!)}
\end{matrix}
\right\},
\]
when varying over $h_1,\ldots, h_n\in F[X]^{<3}$, independent of the masks $\msk_0,\ldots,\msk_{r-1}$. 
\ulrich{Need to write separate relations for each round of WHIR}
\end{lem}

\begin{proof}
First for a single WHIR round, induction on its sumcheck rounds.
$\Box$
\end{proof}